\PassOptionsToPackage{colorlinks=true,allcolors=black}{hyperref}
% REQ: The above line, \PassOptionsToPackage{} must be very first line.

\documentclass[11pt]{article}

% Encoding and font setup for cross-platform reproducibility
\usepackage[T1]{fontenc}
\usepackage{lmodern}

% Improve readability for dense theoretical material
\usepackage{setspace}
\onehalfspacing

% Mathematical notation and table formatting for formal semantics
\usepackage{amsmath, amssymb}
\usepackage{array}
\usepackage[table]{xcolor}
\usepackage{colortbl}
\usepackage{booktabs}
\usepackage{graphicx}

% Theorem environments
\usepackage{amsthm}

% =============================================================================
% Numbering convention:
% - All theorem-like statements share one counter, numbered by section.
% - Prose cites statements by stable names and labels, not by number.
% - External documents cite by stable tag (e.g., se110.def.Substrate).
%   Numbers move when text is edited; names do not.
% =============================================================================
\theoremstyle{plain}
\newtheorem{theorem}{Theorem}[section]
\newtheorem{lemma}[theorem]{Lemma}
\newtheorem{proposition}[theorem]{Proposition}
\newtheorem{constraint}[theorem]{Constraint}
\newtheorem{claim}[theorem]{Claim}
\newtheorem{assumption}[theorem]{Assumption}

\theoremstyle{definition}
\newtheorem{definition}[theorem]{Definition}
\newtheorem{example}[theorem]{Example}

\theoremstyle{remark}
\newtheorem{remark}[theorem]{Remark}
\newtheorem{note}[theorem]{Note}

% Unnumbered versions for front-matter previews/callouts
\newtheorem*{constraint*}{Constraint}
\newtheorem*{definition*}{Definition}
\newtheorem*{example*}{Example}
\newtheorem*{remark*}{Remark}
\newtheorem*{note*}{Note}

% Bibliography management and hyperlink support
\usepackage{url}
\usepackage{natbib}
\bibpunct{(}{)}{;}{a}{}{ }

% Core hyperlink and bookmark infrastructure
\usepackage{hyperref}
\usepackage{bookmark}

% Support multiple author affiliations
\usepackage{authblk}

% Improve typographic quality and line breaking
\usepackage{microtype}

% Callout boxes for examples and emphasis
\usepackage[most]{tcolorbox}

% Compact list formatting for dense conceptual content
\usepackage{enumitem}
\setlist[itemize]{itemsep=0.2em, topsep=0.2em, parsep=0em, partopsep=0em}
\setlist[enumerate]{itemsep=0.2em, topsep=0.2em, parsep=0em, partopsep=0em}

% Paragraph spacing prioritizes readability over traditional indentation
\setlength{\parskip}{0.75em}
\setlength{\parindent}{0em}

% Keywords macro for structured abstract metadata
\providecommand{\keywords}[1]{\vspace{0.5em}\noindent\textbf{Keywords: }\small #1}

% =============================================================================
% Stable named-reference helpers.
% Use labels such as se110.def.Substrate and cite by stable name:
% \DefRef{se110.def.Substrate}{Substrate}
% =============================================================================
\newcommand{\DefRef}[2]{\hyperref[#1]{Definition (#2)}}
\newcommand{\AssumpRef}[2]{\hyperref[#1]{Assumption (#2)}}
\newcommand{\ConstRef}[2]{\hyperref[#1]{Constraint (#2)}}
\newcommand{\ExRef}[2]{\hyperref[#1]{Example (#2)}}
\newcommand{\RemRef}[2]{\hyperref[#1]{Remark (#2)}}
\newcommand{\NoteRef}[2]{\hyperref[#1]{Note (#2)}}
\newcommand{\PropRef}[2]{\hyperref[#1]{Proposition (#2)}}
\newcommand{\SecRef}[1]{Section~\ref{#1}}

% =============================================================================
% Formal notation
% =============================================================================
\newcommand{\Substrate}{\mathcal{S}}
\newcommand{\SubstrateRef}{\mathcal{S}_{\mathrm{ref}}}
\newcommand{\Frameworks}{\mathbb{F}}
\newcommand{\Framework}{F}
\newcommand{\FrameworkOne}{F_1}
\newcommand{\FrameworkTwo}{F_2}
\newcommand{\entails}{\vdash}
\newcommand{\notentails}{\nvdash}
\newcommand{\Asserts}{\mathsf{Asserts}}
\newcommand{\pred}[1]{\mathit{#1}}

% Reusable callout box
\newcommand{\FigureCallout}[2]{%
  \begin{tcolorbox}[
    colback=gray!5,
    colframe=black!40,
    title={#1},
    fonttitle=\bfseries,
    arc=3pt,
    boxrule=0.5pt,
    width=\linewidth,
    enhanced,
    breakable
  ]
  #2
  \end{tcolorbox}
}

% =============================================================================
% PREFERRED CITATION SPECIFICATION
%
% @techreport{case2026common,
%   author      = {Case, Denise M.},
%   title       = {Referential Common Ground: Conditions for Compatibility Between a Bounded Referential Fragment and Interpretive Frameworks},
%   institution = {Preprint},
%   year        = {2026},
%   doi         = {10.48550/arXiv.2601.14271},
%   url         = {https://doi.org/10.48550/arXiv.2601.14271}
% }
% =============================================================================

% ============================================================
\title{
Referential Common Ground:\\
Conditions for Compatibility Between a Bounded Referential Fragment\\
and Interpretive Frameworks
}
% ============================================================

\author{Denise M. Case}
\affil{\small Department of Computer Science and Information Systems\\
Northwest Missouri State University, Maryville, MO, USA\\}
\date{}

\begin{document}
\maketitle
\vspace{-2em}

\begin{abstract}
  Shared records must often serve parties who interpret them differently.
  The neutral-substrate result shows that a foundational layer limited to
  referential commitments and permitted attributions is neutral relative to
  every admissible framework, provided the referential commitments are jointly
  consistent with each admissible framework.
  That proviso, Referential Common Ground, was left as an assumption for each
  realization to justify.

  This paper asks what can be proved about it.
  Admissibility alone does not supply it: when the admissibility conditions
  beyond consistency place no restriction on a framework's content, a
  commitment set is compatible with every admissible framework exactly when
  every valuation satisfies it, so no nontrivial referential fragment
  qualifies.
  The needed condition must therefore constrain frameworks.
  We state it without reference to any fragment's commitments: a framework is
  neutral on a vocabulary when it permits every valuation of that vocabulary.
  A framework is neutral on a vocabulary exactly when it is jointly satisfiable
  with every satisfiable theory over it; every framework is either neutral or
  conflicts with some satisfiable theory over the vocabulary; and a satisfiable
  framework written entirely outside the vocabulary is neutral.

  The results are proved for classical propositional logic and checked in
  Lean 4.
  They are relative to a stated fragment and framework class.
  They do not establish common ground for the full class of admissible
  frameworks, and they do not justify that a fragment's commitments are
  referential.
\end{abstract}


\keywords{
  Formal ontology;
  neutral substrates;
  referential common ground;
  interpretive frameworks;
  framework admissibility;
  joint consistency;
  persistent interpretive disagreement
}

\begin{note*}[Citation Convention]
  Definitions, assumptions, constraints, and examples in this paper carry stable
  names (for example, \texttt{se110.def.NeutralOn}).
  This paper cites statements by stable name rather than by number,
  so citations remain valid even if numbering changes.
\end{note*}

% =============================================================================
% Preview of Main Result
% =============================================================================

\bigskip
\noindent\textbf{Preview of Main Result.}
\emph{
  Let $V$ be the vocabulary of a satisfiable referential fragment.
  A framework is neutral on $V$ if it permits every valuation of $V$.
  A framework is neutral on $V$ exactly when it is jointly satisfiable with
  every satisfiable theory over $V$, so a neutral framework is compatible with
  the fragment whatever its content.
  Admissibility alone does not give this: when it imposes no restriction beyond
  consistency, only valid commitment sets are compatible with every admissible
  framework.
}

\medskip
\noindent
\emph{Informally.} A framework that says nothing about what a record is about
cannot contradict what the record is about.
A framework that does say something can.
Common ground for a fragment therefore holds for the frameworks that stay
silent on the fragment's vocabulary.
For the others it is a question about the framework's content.

% =============================================================================
\section{Introduction}
\label{sec:introduction}
% =============================================================================


The purpose of this paper is to determine what can be proved about Referential
Common Ground, the assumption on which the neutral-substrate result
\citep{case2026neutral} rests.

The central claim is:

\begin{quote}
  Referential Common Ground cannot be derived from admissibility alone.
  For a bounded referential fragment it can be established relative to a class
  of frameworks characterized by what they permit on the fragment's
  vocabulary, without reference to the fragment's commitments.
  A framework in that class is compatible with every satisfiable theory over
  the vocabulary, and every other framework conflicts with some such theory.
\end{quote}

The mechanism is the \emph{interface} of a framework on a vocabulary: the set of
valuations of the vocabulary that the framework permits.
Compatibility of a theory over the vocabulary with the framework depends only
on that interface.
A framework that permits every valuation cannot contradict any satisfiable
theory over the vocabulary.

The paper proceeds as follows.
\SecRef{sec:related-work} credits the foundations and states the contribution.
\SecRef{sec:problem-setting} states the problem.
\SecRef{sec:fragments} defines fragments.
\SecRef{sec:independence} shows that admissibility alone does not give common ground.
\SecRef{sec:interface} defines the interface a framework permits and the scope condition.
\SecRef{sec:characterization} characterizes the scope condition and proves the dichotomy.
\SecRef{sec:example} gives a worked example.
\SecRef{sec:limits} states the limits of this work.
\SecRef{sec:discussion} discusses the practical significance and limits of the work.
\SecRef{sec:conclusion} concludes.

% =============================================================================
\section{Related Work and Contribution}
\label{sec:related-work}
% =============================================================================

The neutral-substrate result \citep{case2026neutral} restricts a shared
foundational layer to referential commitments and permitted attributions, and
shows that this restriction yields neutrality relative to every admissible
framework under Referential Common Ground.
Its representational background (reification, provenance, context logic,
argumentation, standpoint logic) is given there and is not repeated.
This paper takes Referential Common Ground itself as its subject.

The logical facts used here are classical.
Joint consistency of satisfiable propositional theories that share no atom is
the simplest instance of the phenomenon behind Robinson's joint consistency
theorem \citep{robinson1956}, which is closely related to Craig interpolation
\citep{craig1957}.
In first-order logic with equality the disjoint-vocabulary case fails in
general, because a theory can constrain the size of the domain.
This is why the paper works propositionally.

The scope condition is close in spirit to conservative extensions as used in
ontology modularity \citep{cuencagrau2008modular}.
Remark~\ref{rem:conservative} states the connection.
The purpose differs: here the condition characterizes frameworks that cannot
contradict a referential fragment, in a setting where frameworks are admitted by
documented provenance and not by logical form.

\textbf{The contribution.}
Each ingredient is known.
The contribution is the formulation.
It gives a condition on frameworks that is stated independently of the
fragment, proves that this condition is exactly what is needed uniformly over
fragments on a vocabulary, and proves that admissibility alone cannot supply
any such guarantee.


% =============================================================================
\section{Problem Setting}
\label{sec:problem-setting}
% =============================================================================

This paper builds on the neutral-substrate result \citep{case2026neutral}.
That result uses the following notions, restated here only as far as needed.
A \emph{commitment theory} is a set of object-language propositions.
A consequence relation $\entails$ relates theories to propositions.
It has three structural properties: every member of a theory is entailed by it,
generalized cut, and a proposition together with its negation entails the
contradiction $\bot$.
A theory is \emph{consistent} when it does not entail $\bot$.
A \emph{framework} $\Framework$ is an interpretive position with a commitment
theory, and we write $\Framework$ also for that theory.
A framework is \emph{admissible} when its theory is consistent, it is
evidentially grounded, and it has a documented interpretive function.
A commitment set $T$ is \emph{framework-compatible} when
$T \cup \Framework \notentails \bot$ for every admissible framework
$\Framework$.
\emph{Referential Common Ground} is the assumption that the referential
commitments $\SubstrateRef$ of a substrate are framework-compatible.

The neutral-substrate result leaves Referential Common Ground as an assumption
that a realization must justify.
Three ways of dispensing with it are unavailable.
It cannot be derived from admissibility, as \SecRef{sec:independence} shows.
It cannot be secured by building agreement with $\SubstrateRef$ into the
definition of admissibility, because the result would then hold by definition
and say nothing about frameworks.
It cannot be asserted for the full class of admissible frameworks without an
argument about every framework that could be admitted.

What can be done is narrower.
A legitimate result must satisfy five requirements.
\begin{enumerate}
  \item It is relative to a stated fragment and a stated class of frameworks,
        and does not claim the full class.
  \item The fragment is nontrivial: a fragment that rules nothing out is
        compatible with every satisfiable framework, and a result about it
        says nothing.
  \item The class of frameworks is characterized without reference to the
        fragment's commitments.
  \item A constructive compatibility theorem and a countermodel showing why a
        proposed condition is inadequate are both acceptable outcomes.
  \item That the fragment's commitments are genuinely referential is a premise
        supplied from outside this paper.
\end{enumerate}



% =============================================================================
\section{Fragments}
\label{sec:fragments}
% =============================================================================

All results are stated for classical propositional logic, the concrete
instance of the consequence relation of \SecRef{sec:problem-setting} in which
the formal development is carried out.

Let $\mathsf{At}$ be a set of atoms.
Formulas are built from atoms, the contradiction $\bot$, negation, and
conjunction.
A \emph{valuation} $v$ assigns truth or falsity to each atom and extends to all
formulas in the usual way.
A valuation \emph{satisfies} a theory $T$, written $v \models T$, when it makes
every member of $T$ true.
A theory is \emph{satisfiable} when some valuation satisfies it.
A theory \emph{semantically entails} a formula when every valuation that
satisfies the theory makes the formula true.
This relation has the three structural properties required of $\entails$, and a
theory is consistent exactly when it is satisfiable.
Satisfiability therefore stands in for consistency throughout.

A \emph{vocabulary} is a set $V \subseteq \mathsf{At}$.
A formula \emph{lies in} $V$ when every atom occurring in it belongs to $V$,
and a theory lies in $V$ when all of its members do.

\begin{lemma}[Agreement]
  If two valuations agree on every atom of $V$, they satisfy the same theories
  that lie in $V$.
\end{lemma}

\begin{proof}
  By induction on formulas, a formula lying in $V$ is true under a valuation
  exactly when it is true under any valuation that agrees with it on $V$.
\end{proof}

\begin{proposition}[Joint Consistency over Disjoint Vocabularies]
  \label{se110.proposition.SatisfiableUnionOfDisjoint}
  Let $T$ lie in $V$ and $U$ lie in $W$, where $V$ and $W$ have no atom in
  common.
  If $T$ and $U$ are each satisfiable, then $T \cup U$ is satisfiable.
\end{proposition}

\begin{proof}
  Let $v_1 \models T$ and $v_2 \models U$.
  Let $v$ agree with $v_1$ on $V$ and with $v_2$ outside $V$.
  By agreement $v \models T$.
  Since $W$ is disjoint from $V$, $v$ agrees with $v_2$ on $W$, so
  $v \models U$.
\end{proof}

\begin{definition}[Referential Fragment]
  \label{se110.def.ReferentialFragment}
  A \emph{referential fragment} is a vocabulary $V$ together with a theory $R$
  lying in $V$ such that
  \begin{enumerate}
    \item $R$ is satisfiable;
    \item $R$ is nontrivial: some valuation does not satisfy $R$; and
    \item $V$ is bounded: it is contained in a finite set of atoms.
  \end{enumerate}
\end{definition}

Condition~1 keeps the fragment consistent.
Condition~2 excludes the fragments for which the question is vacuous
(\PropRef{se110.proposition.CompatibleBareIffValid}{Bare Compatibility}).
Condition~3 records that the fragments of interest are small.
No result below uses it.

The definition fixes the logical shape of a fragment.
It does not say that the commitments of $R$ are referential.
That they are fixed by a referential regime, and that the regime is justified
for the intended referents, is a premise supplied by the realization, as in
\citet{case2026neutral}.

% =============================================================================
\section{Independence}
\label{sec:independence}
% =============================================================================

To see what admissibility contributes on its own, consider the setting in which
the evidential-grounding and documentation conditions hold for every framework.
A framework is then admissible exactly when its theory is consistent, that is,
satisfiable.
We call this the \emph{bare setting}.
It is the setting in which admissibility imposes no restriction on a
framework's content beyond consistency: the grounding and documentation
conditions concern provenance and presentation, not content.

\begin{proposition}[Bare Compatibility]
  \label{se110.proposition.CompatibleBareIffValid}
  In the bare setting, a commitment set $T$ is framework-compatible if and only
  if every valuation satisfies $T$.
  Consequently no referential fragment is framework-compatible.
\end{proposition}

\begin{proof}
  Suppose every valuation satisfies $T$, and let $\Framework$ be admissible.
  Any valuation satisfying $\Framework$ satisfies $T \cup \Framework$, so the
  union is satisfiable.

  Conversely, suppose some valuation $v$ does not satisfy $T$.
  Then some $\varphi \in T$ is false under $v$.
  The theory $\{\neg\varphi\}$ is satisfied by $v$, so it is admissible in the
  bare setting.
  Any valuation satisfying $T$ makes $\varphi$ true, so none satisfies
  $T \cup \{\neg\varphi\}$.
  Thus $T$ is not framework-compatible.

  A referential fragment is nontrivial, so its commitments are not satisfied
  by every valuation.
\end{proof}

The proof is also the countermodel promised in \SecRef{sec:problem-setting}.
For any commitment of a nontrivial fragment, the single-member theory that
denies it is satisfiable and therefore admissible in the bare setting, and it
conflicts with the fragment.
This does not contradict the neutral-substrate result, where Referential Common
Ground is an assumption.
It shows that the assumption cannot be replaced by a consequence of
admissibility.


% =============================================================================
\section{Interface and Scope}
\label{sec:interface}
% =============================================================================

Fix a vocabulary $V$.

\begin{definition}[Interface Permission]
  \label{se110.def.Permits}
  A framework $\Framework$ \emph{permits} a valuation $u$ on $V$ when some
  valuation satisfying $\Framework$ agrees with $u$ on every atom of $V$.
\end{definition}

The valuations of $V$ that $\Framework$ permits form its \emph{interface} on
$V$.
The next result shows that compatibility with a theory over $V$ depends only on
the interface.

\begin{proposition}[Compatibility Reduces to the Interface]
  \label{se110.proposition.SatisfiableUnionIffPermits}
  Let $R$ lie in $V$.
  Then $R \cup \Framework$ is satisfiable if and only if $\Framework$ permits
  some valuation that satisfies $R$.
\end{proposition}

\begin{proof}
  If $w$ satisfies $R \cup \Framework$, then $w$ satisfies $R$, and
  $\Framework$ permits $w$ because $w$ satisfies $\Framework$ and agrees with
  itself.
  Conversely, let $u \models R$ and let $w \models \Framework$ agree with $u$ on
  $V$.
  Since $R$ lies in $V$, agreement gives $w \models R$.
  Hence $w \models R \cup \Framework$.
\end{proof}

\begin{definition}[Neutral on a Vocabulary]
  \label{se110.def.NeutralOn}
  A framework $\Framework$ is \emph{neutral on} $V$ when it permits every
  valuation of $V$.
\end{definition}

The condition mentions only $V$ and $\Framework$.
It does not mention any fragment's commitments.
A neutral framework is satisfiable.

\begin{definition}[Avoiding a Vocabulary]
  \label{se110.def.AvoidsVocabulary}
  A framework $\Framework$ \emph{avoids} $V$ when it lies in some vocabulary
  $W$ with no atom in common with $V$.
\end{definition}

\begin{proposition}[Avoidance Gives Neutrality]
  \label{se110.proposition.NeutralOnOfAvoids}
  A satisfiable framework that avoids $V$ is neutral on $V$.
\end{proposition}

\begin{proof}
  Let $\Framework$ lie in $W$, disjoint from $V$, and let $w_2 \models
  \Framework$.
  Given a valuation $u$ of $V$, let $w$ agree with $u$ on $V$ and with $w_2$
  outside $V$.
  Since $W$ is disjoint from $V$, $w$ agrees with $w_2$ on $W$, so
  $w \models \Framework$ by agreement.
  Thus $\Framework$ permits $u$.
\end{proof}

Avoidance can be checked by reading the framework's vocabulary.
Neutrality is semantic and weaker: a framework may mention atoms of $V$
vacuously, for example through $p \vee \neg p$, and still be neutral.

% =============================================================================
\section{Characterization and Dichotomy}
\label{sec:characterization}
% =============================================================================

\begin{proposition}[Neutrality Characterization]
  \label{se110.proposition.NeutralOnIffCompatible}
  A framework $\Framework$ is neutral on $V$ if and only if $R \cup \Framework$
  is satisfiable for every satisfiable theory $R$ lying in $V$.
\end{proposition}

\begin{proof}
  Suppose $\Framework$ is neutral on $V$ and $R \models$-satisfiable by some
  $u$.
  Since $\Framework$ permits $u$,
  \PropRef{se110.proposition.SatisfiableUnionIffPermits}{Compatibility Reduces
    to the Interface} gives that $R \cup \Framework$ is satisfiable.

  Conversely, suppose $\Framework$ is compatible with every satisfiable theory
  lying in $V$, and let $u$ be a valuation of $V$.
  Let $D_u$ contain $p$ for each atom $p \in V$ that $u$ makes true and
  $\neg p$ for each atom $p \in V$ that $u$ makes false.
  Then $D_u$ lies in $V$ and is satisfied by $u$.
  By assumption some $w$ satisfies $D_u \cup \Framework$.
  Then $w \models \Framework$, and $w$ agrees with $u$ on $V$ because it
  satisfies $D_u$.
  So $\Framework$ permits $u$.
\end{proof}

\begin{proposition}[Scope Dichotomy]
  \label{se110.proposition.Dichotomy}
  For every framework $\Framework$ and vocabulary $V$, either $\Framework$ is
  neutral on $V$, or some satisfiable theory $R$ lying in $V$ is such that
  $R \cup \Framework$ is not satisfiable.
\end{proposition}

\begin{proof}
  Classically, either $\Framework$ is neutral on $V$ or it is not.
  If it is not, \PropRef{se110.proposition.NeutralOnIffCompatible}{Neutrality
    Characterization} yields a satisfiable theory lying in $V$ that is not
  jointly satisfiable with $\Framework$.
\end{proof}

The second case of the dichotomy supplies a witness: a concrete satisfiable
theory over $V$ that the framework cannot coexist with.
Whether that theory is the commitment set of a given fragment is a question
about the fragment.
For a fixed fragment, weaker conditions than neutrality can suffice, as
\PropRef{se110.proposition.SatisfiableUnionIffPermits}{Compatibility Reduces
  to the Interface} shows.
Neutrality is the weakest condition on a framework alone that gives
compatibility with every fragment on $V$.

\begin{definition}[Fragment-Relative Common Ground]
  \label{se110.def.FragmentCommonGround}
  Let $(V,R)$ be a referential fragment and $\mathcal{K}$ a class of frameworks.
  The fragment has \emph{common ground relative to} $\mathcal{K}$ when
  $R \cup \Framework$ is satisfiable for every $\Framework \in \mathcal{K}$.
\end{definition}

This is not Referential Common Ground as used in the neutral-substrate result.
When $\mathcal{K}$ is the class of all admissible frameworks and $R$ is the
substrate's referential commitments, it agrees with the first condition of that
assumption.
This paper establishes it only for classes of neutral frameworks.

\begin{proposition}[Common Ground from Neutrality]
  \label{se110.proposition.FragmentCommonGroundOfNeutral}
  Let $(V,R)$ be a referential fragment.
  If every framework in $\mathcal{K}$ is neutral on $V$, then the fragment has
  common ground relative to $\mathcal{K}$.
\end{proposition}

\begin{proof}
  $R$ is satisfiable and lies in $V$.
  Apply \PropRef{se110.proposition.NeutralOnIffCompatible}{Neutrality
    Characterization} to each $\Framework \in \mathcal{K}$.
\end{proof}

\begin{remark}[Relation to Conservative Extensions]
  \label{rem:conservative}
  A framework $\Framework$ is neutral on $V$ exactly when every formula lying
  in $V$ that $\Framework$ entails is valid.
  The forward direction is immediate.
  The converse uses compactness of propositional logic: if $\Framework$ forbids
  $u$, some finite part of $D_u$ already conflicts with $\Framework$, so
  $\Framework$ entails the negation of that finite conjunction, a formula
  lying in $V$ that $u$ falsifies.
  In the terminology of ontology modularity \citep{cuencagrau2008modular}, this
  says that $\Framework$ is a conservative extension of the empty theory with
  respect to $V$.
  This remark is not part of the formalization.
\end{remark}

% =============================================================================
\section{Worked Example}
\label{sec:example}
% =============================================================================

\begin{example}[Decision Record]
  \label{se110.example.DecisionRecord}
  Consider a decision record with four atoms:
  $\pred{decisionMade}$ (a decision was made),
  $\pred{occurredAt}$ (it occurred at the recorded time),
  $\pred{causedByScore}$ (the denial was caused by the score), and
  $\pred{justifiedByPolicy}$ (the cited policy justified the decision).
  Let $V = \{\pred{decisionMade}, \pred{occurredAt}\}$ and
  $R = \{\pred{decisionMade}, \pred{occurredAt}\}$.
  The theory $R$ is satisfied by the valuation that makes every atom true, and
  not satisfied by the valuation that makes every atom false.
  So $(V,R)$ is a referential fragment.
\end{example}

The example is conditional: it presupposes a realization that justifies
$\pred{decisionMade}$ and $\pred{occurredAt}$ as referential commitments.
The atoms do not by themselves establish that.

\textbf{A neutral framework.}
Let $\Framework_1 = \{\pred{causedByScore}\}$, a causal claim.
It lies in $\{\pred{causedByScore}, \pred{justifiedByPolicy}\}$, which is
disjoint from $V$, and it is satisfiable.
By \PropRef{se110.proposition.NeutralOnOfAvoids}{Avoidance Gives Neutrality}
it is neutral on $V$.
By \PropRef{se110.proposition.NeutralOnIffCompatible}{Neutrality
  Characterization}, $R \cup \Framework_1$ is satisfiable.
The fragment has common ground relative to any class of frameworks of this kind.

\textbf{A conflicting framework.}
Let $\Framework_2 = \{\neg\pred{occurredAt}\}$, a framework that disputes the
recorded time.
It is satisfied by the valuation that makes every atom false, so it is
admissible in the bare setting.
Every valuation satisfying $R$ makes $\pred{occurredAt}$ true, so no valuation
satisfies $R \cup \Framework_2$.
By the dichotomy, $\Framework_2$ is not neutral on $V$, and $R$ itself is the
witness.
By \PropRef{se110.proposition.CompatibleBareIffValid}{Bare Compatibility}, the
fragment is not framework-compatible in the bare setting, and $\Framework_2$
is the framework that shows it.

\textbf{Reading.}
The conflict is not resolved by choosing a side.
In the neutral-substrate construction, $\Framework_2$'s dispute with the
recorded time belongs in the record as an attributed assertion.
The commitment to the recorded time is then not common ground for
$\Framework_2$, and a class of frameworks for which it is common ground must
exclude $\Framework_2$.


% =============================================================================
\section{Limits}
\label{sec:limits}
% =============================================================================

\textbf{This paper does not establish Referential Common Ground for the full
  class of admissible frameworks.}
Its results are relative to a stated fragment and a stated class of
frameworks.

\textbf{This paper does not justify referential fixing.}
That the commitments of a fragment are genuinely referential is a premise of the
realization.
Nothing here derives it.

\textbf{This paper does not show that realistic frameworks are neutral.}
It states the condition a class of frameworks must meet and proves what the
condition delivers.
Whether a deployed class meets it is a question about that class.

\textbf{This paper does not treat identity-bearing fragments.}
Fragments whose commitments assert identity, co-reference, or persistence are
out of scope.
In first-order logic with equality, theories over disjoint vocabularies can
conflict through the size of the domain, so the disjointness argument does not
extend.

\textbf{This paper does not provide a general compatibility algorithm.}
Avoidance is checked by reading the framework's vocabulary.
Neutrality is a semantic condition and is not claimed to be decidable in
general.

\textbf{This paper does not claim more than classical propositional logic.}
Results that depend on classical features hold for that logic, not for every
consequence relation with the three structural properties.

% =============================================================================
\section{Discussion}
\label{sec:discussion}
% =============================================================================

Referential Common Ground, as used in the neutral-substrate result, quantifies
over the whole class of admissible frameworks, including frameworks not known
when the substrate is designed.
The results here do not weaken that quantifier.
They show that, for a fragment, the quantifier can be discharged for exactly
the frameworks that are neutral on the fragment's vocabulary, and that no weaker
condition on a framework alone works uniformly over fragments on that
vocabulary.

The practical reading has three parts.
First, avoidance can be checked by reading a framework's vocabulary against the
fragment's.
A framework that avoids the fragment's vocabulary and is satisfiable is neutral.
Second, a framework that does use the fragment's vocabulary is a boundary case,
not an excluded one.
Whether it is compatible is a question about its content.
If it disputes a fragment commitment, the dispute belongs in the record as an
attributed assertion, and that commitment is not common ground for the
framework.
Third, a deployment could require neutrality on the fragment's vocabulary as a
condition of admission.
That is a design choice and not a consequence of the results.
It would make the admissibility gate depend on the fragment, and the
neutral-substrate result already observes that the gate is not neutral.

The propositional setting is a limitation as well as a convenience.
It gives clean proofs and excludes the identity-bearing fragments for which the
question matters most.
Extending the analysis to those fragments needs identity criteria and is
outside this paper.


% =============================================================================
\section{Conclusion}
\label{sec:conclusion}
% =============================================================================

Referential Common Ground cannot be obtained from admissibility alone: when
admissibility imposes no restriction on content beyond consistency, only valid
commitment sets are framework-compatible, and no nontrivial referential
fragment qualifies.

The condition that does the work constrains frameworks.
A framework is neutral on a vocabulary when it permits every valuation of that
vocabulary, a condition stated without reference to any fragment's commitments.
Neutral frameworks are exactly those compatible with every satisfiable theory
over the vocabulary.
Every other framework conflicts with some such theory.

The result is conditional and small.
It is relative to a stated fragment and framework class, it makes no claim about
the full class of admissible frameworks, and it does not justify that a
fragment's commitments are referential.
It turns an unexamined assumption into a precise obligation: to show that the
frameworks in a deployed class are neutral on the vocabulary, or to say which
are not.


% =============================================================================
\section*{Statements and Declarations}
% =============================================================================

\subsection*{Acknowledgments} The author thanks the anonymous reviewers
of an earlier version of this work for comments that helped improve
the framing, organization, and presentation of the work.

\subsection*{Author Contributions}
The author is the sole contributor to this work and
is responsible for all aspects of the research, authorship, and publication.

\subsection*{Use of AI-Assisted Tools}
AI-assisted tools were used for editing, formatting, and consistency checking.
The author reviewed all suggestions and is solely responsible for the content.

\subsection*{Declaration of Conflicting Interest}
The author declares no potential conflicts of interest.

\begin{thebibliography}{99}

  \bibitem[Case(2026)]{case2026neutral}
  Case, D.~M. (2026).
  Neutral substrates: A design constraint for shared records under persistent
  interpretive disagreement.
  Preprint, arXiv:2601.14271.

  \bibitem[Craig(1957)]{craig1957}
  Craig, W. (1957).
  Three uses of the Herbrand--Gentzen theorem in relating model theory and
  proof theory.
  \emph{The Journal of Symbolic Logic}, 22(3):269--285.

  \bibitem[Cuenca~Grau et~al.(2008)]{cuencagrau2008modular}
  Cuenca~Grau, B., Horrocks, I., Kazakov, Y., and Sattler, U. (2008).
  Modular reuse of ontologies: Theory and practice.
  \emph{Journal of Artificial Intelligence Research}, 31:273--318.

  \bibitem[Robinson(1956)]{robinson1956}
  Robinson, A. (1956).
  A result on consistency and its application to the theory of definition.
  \emph{Koninklijke Nederlandse Akademie van Wetenschappen, Proceedings,
    Series A}, 59:47--58; also \emph{Indagationes Mathematicae}, 18:47--58.

\end{thebibliography}

\end{document}
